\section{Q\&A 与实践提醒}

\subsection{现场问答要点}

在 Q\&A 环节，Chelsea Finn 强调“observe → instrument → govern”三步，特别是在 emergent drift 发生后要找出具体 prompt/data/attention artifact。

\begin{table}[H]
\centering
\begin{tabular}{p{4cm}p{5cm}}
\toprule
问题主题 & 核心回答 \\
\midrule
Emergent stability & 多指标对齐后才能宣布 emergence 成熟 \\
Prompt drift & insist on prompt log + rollback plan \\
Tooling stack & attention/plan/metric traces 必须 realtime + archival \\
\bottomrule
\end{tabular}
\caption{Q\&A 环节精华。}
\end{table}

\begin{knowledgebox}{Q\&A 记忆点}
把 Q\&A 中的 keyword 写进 `qa.md`，每次 incident review 时对照，确保没有遗漏高风险提醒。
\end{knowledgebox}

\subsection{实践提醒}

\begin{figure}[H]
\centering
\includegraphics[width=0.85\textwidth]{slides-images/slide-41.jpg}
\caption{Q\&A 与实践提醒 slide（slide-41）。}
\end{figure}

\begin{importantbox}{实践落地提醒}
\begin{itemize}
    \item 记录 prompt version + sampling seed；
    \item 每个 emergent drill 都要产出 lessons learned + summary note；
    \item 在 governance panel 会议上 review incident grade-book + prompt log。
\end{itemize}
\end{importantbox}

\subsection{本章小结}

Q\&A 让讲座的 abstract insights 成为 checklist；只要把问答中的 artifact 持续写出来，治理就始终有据可循。

